\section{The Angle}

The weak mixing angle is the sharpest dimensionless number the model predicts, and it must be stated with its real
caveats.

The bare value at the unification scale is three eighths. This is pure group theory from the charges of the sixteen-spinor
of the larger organizing group, with its hypercharge normalization, and it is exact, a number that comes from counting
directions and nothing else.

The low-energy value is a genuine prediction, not a fit. Take only two measured inputs, the electromagnetic coupling and
the strong coupling at the relevant scale, demand that the three couplings meet at one unification scale, and the value
of the weak mixing angle at low energy comes out as an output. The result, exactly as the computation gives it: with
supersymmetric-like running it is about 0.231, matching the measured 0.2312, and with bare Standard Model running it is
about 0.207, the known small miss. The running and the wrong-charge control are in the open codebase \cite{cluesurf2026repo}.

So the headline match to the measured angle assumes supersymmetric-like content in the running. The number could have
come out wrong, since only the two couplings are fed in and the angle is an
output, and a wrong hypercharge normalization predicts a clearly wrong angle, the discriminating control. So the
structure is doing the work, the three-eighths is forced by the charge organization, and the prediction carries its
running caveat with it.
